\chapter{AIXI and Gödel Machines}

\section{Introduction}

The preceding chapters developed discrete dynamical systems on finite state spaces and explored Kolmogorov complexity as a measure of description length. We now turn to two landmark theoretical constructions that connect these ideas to the problem of \textit{optimal sequential decision-making}: Hutter's AIXI and Schmidhuber's Gödel Machine.

Both constructions grapple with the same fundamental question: what does it mean for an agent to act optimally in an unknown environment? AIXI answers this by maintaining a Bayesian mixture over all computable environments, weighted by a universal prior. The Gödel Machine answers it by searching for provably beneficial self-modifications to its own code. Both are incomputable in their full generality, yet both illuminate the structure of the problem in ways that computable approximations cannot.

The thread connecting them to our framework is direct. In the notation of Chapter 1, a discrete dynamical system is a pair $(X, f)$ where $f: X \to X$. When we allow the map itself to change---$(f, x) \mapsto (\phi(f, x),\, f(x))$---we enter the territory of self-modifying systems. AIXI reasons \textit{about} an unknown $f$ without modifying itself. Gödel Machines modify their own $f$ subject to a proof constraint. The general $(f, x) \mapsto (\phi(f,x), f(x))$ framework studied in this text imposes no such constraint, allowing us to study the full landscape of self-modification dynamics.

We begin with the foundation on which AIXI rests: Solomonoff's theory of inductive inference.



\section{Solomonoff Induction}

\subsection{13.2.1 The Prediction Problem}

Consider the following problem. You observe a binary sequence $x_1, x_2, \ldots, x_n$ and wish to predict $x_{n+1}$. You do not know the process that generated the sequence. It might be a deterministic rule, a biased coin, a pseudorandom generator, or any other mechanism.

Classical statistics requires you to choose a model class (e.g., "the data are i.i.d. Bernoulli") before you can do inference. If the true source lies outside your model class, your predictions may be systematically wrong. Solomonoff induction eliminates this dependence on a pre-chosen model class by considering \textit{all} computable hypotheses simultaneously.

\subsection{13.2.2 The Universal Prior}

Fix a prefix-free universal Turing machine $U$. (Prefix-free means that the set of programs on which $U$ halts forms a prefix-free set: no halting program is a proper prefix of another. This is needed to make the prior a proper probability measure.)

\begin{definition}[Solomonoff's Universal Prior]
For a finite binary string $x \in \{ 0,1\} ^\textit{$, the }universal a priori probability\textit{ (or }Solomonoff prior*) is

\[M(x) = \sum_{\substack{p \in \{ 0,1\} ^\textit{ \\ U(p) = x^}}} 2^{-|p|}\]

where the sum is over all programs $p$ such that $U$ with input $p$ outputs a string starting with $x$ (written $x^*$ to denote that $x$ is a prefix of the output), and $|p|$ is the length of $p$ in bits.

Equivalently, $M(x)$ is the probability that the universal machine $U$, fed with a sequence of fair coin flips as input, produces an output beginning with $x$.
\end{definition}

\begin{example}
Consider three binary strings of length 20:
- $a = 00000000000000000000$ (all zeros)
- $b = 01010101010101010101$ (alternating)
- $c = 10110100011010001101$ (apparently random)

Under the Solomonoff prior, $M(a) > M(b) > M(c)$ (up to constant factors). The string $a$ can be generated by very short programs ("print 0 twenty times"), so many short programs contribute to the sum. The string $b$ requires a slightly longer program. The string $c$, if truly incompressible, can essentially only be produced by a program that contains it as a literal, so the shortest contributing program has length $\approx 20$ bits, giving $M(c) \approx 2^{-20}$.
\end{example}

\subsection{13.2.3 The Dominance Property}

The crucial property of $M$ is that it is \textit{universal} among computable probability measures, in the following precise sense.

\begin{theorem}[Dominance, Solomonoff 1964, Zvonkin and Levin 1970]
\textit{For any computable probability measure $\mu$ on $\{ 0,1\} ^}$, there exists a constant $c_\mu > 0$ (depending on $\mu$ but not on $x$) such that for all $x \in \{ 0,1\} ^\textit{$:}

\[M(x) \geq c_\mu \cdot \mu(x).\]

\textit{Proof sketch.} Since $\mu$ is computable, there exists a program $p_\mu$ that, given random bits as input, samples from $\mu$. This program has some fixed length $|p_\mu|$. The contribution of programs that begin with $p_\mu$ to the sum defining $M(x)$ is at least $2^{-|p_\mu|} \cdot \mu(x)$. Setting $c_\mu = 2^{-|p_\mu|}$ gives the result. 
\end{proof}

This theorem says that $M$ assigns probability at least as large (up to a multiplicative constant) as \textit{any} computable measure. It "dominates" the entire class of computable distributions. The constant $c_\mu = 2^{-|p_\mu|}$ is small for complex hypotheses and large for simple ones, which is precisely Occam's razor implemented in a mathematically rigorous way.

\end{theorem}

\subsection{13.2.4 Convergence of Solomonoff Prediction}

Given observations $x_1, \ldots, x_n$ generated by some computable measure $\mu$, the Solomonoff predictor assigns

\[M(x_{n+1} \mid x_1 \ldots x_n) = \frac{M(x_1 \ldots x_n x_{n+1})}{M(x_1 \ldots x_n)}.\]

\begin{theorem}[Solomonoff's Convergence Theorem, Solomonoff 1978]
\textit{If the data are generated by a computable measure $\mu$, then the expected total KL divergence between $\mu$'s predictions and $M$'s predictions is bounded:}

\[\sum_{n=1}^{\infty} \mathbb{E}_\mu \left[ D_{\mathrm{KL}}\!\left(\mu(\cdot \mid x_{<n}) \,\|\, M(\cdot \mid x_{<n})\right) \right] \leq \ln \frac{1}{c_\mu} = |p_\mu| \cdot \ln 2.\]

\textit{In particular, $D_{\mathrm{KL}}(\mu(\cdot \mid x_{<n}) \| M(\cdot \mid x_{<n})) \to 0$ with $\mu$-probability 1.}

This is a remarkable result: the universal predictor \textit{learns any computable environment}, and the total prediction error over the agent's entire lifetime is bounded by the complexity of the true environment. Simple environments are learned quickly; complex ones take longer, but all are learned eventually.
\end{theorem}

\begin{example}
Suppose the true environment produces the sequence $0, 1, 0, 1, 0, 1, \ldots$ deterministically. After observing just a few alternations, the Solomonoff predictor assigns overwhelming probability to programs that generate the alternating pattern (these are short and thus heavily weighted). Programs generating $0,1,0,1,0,1,0,0,...$ (which agree on the first $n$ bits but then diverge) are also present in the mixture, but their collective weight shrinks as more confirming data arrives. The predictor rapidly converges to the correct rule.
\end{example}

\subsection{13.2.5 Incomputability}

\begin{proposition}
\textit{$M(x)$ is not computable.}

\begin{proof}
Computing $M(x)$ requires summing $2^{-|p|}$ over all programs $p$ that produce output beginning with $x$. Determining whether $U(p)$ halts and produces output starting with $x$ is at least as hard as the halting problem. 

However, $M(x)$ is \textit{enumerable from above} (upper semi-computable): we can run all programs in parallel, and as we discover more programs that output strings starting with $x$, our running estimate of $M(x)$ can only increase. We can approximate $M(x)$ from above to arbitrary precision, but we can never know when we have converged.

This incomputability is not a minor technical inconvenience---it is the price of universality. Any computable predictor must fail on some computable environments (this follows from diagonalization), whereas $M$ succeeds on all of them, at the cost of being incomputable.



\end{proof}

\end{proposition}

\section{AIXI}

\subsection{13.3.1 The Agent-Environment Framework}

We now extend from passive prediction to \textit{active} sequential decision-making. Consider an agent interacting with an environment over discrete time steps $t = 1, 2, 3, \ldots, m$ (where $m$ may be finite or a horizon parameter):

1. At time $t$, the agent selects an action $a_t \in \mathcal{A}$.
2. The environment responds with an observation $o_t \in \mathcal{O}$ and a reward $r_t \in \mathcal{R} \subset [0, 1]$.
3. Both agent and environment may depend on the entire interaction history.

The environment is modeled as a computable function (program) $q$ that maps action histories to observation-reward pairs:

\[q: (a_1 o_1 r_1 \ldots a_{t-1} o_{t-1} r_{t-1} a_t) \mapsto (o_t, r_t).\]

The agent's goal is to maximize cumulative reward $\sum_{t=1}^m r_t$.

\subsection{13.3.2 The AIXI Agent}

\begin{definition}[AIXI, Hutter 2000]
The AIXI agent selects actions by maximizing expected future reward, where the expectation is taken over all computable environments weighted by the universal prior:

\[a_t^* = \arg\max_{a_t} \sum_{o_t, r_t} \ldots \max_{a_m} \sum_{o_m, r_m} \left[\sum_{k=t}^{m} r_k \right] \sum_{\substack{q \text{ consistent} \\ \text{with history}}} 2^{-|q|}\]

where "consistent with history" means $q(a_1 o_1 r_1 \ldots a_{k-1} o_{k-1} r_{k-1} a_k) = (o_k, r_k)$ for all $k < t$.

Let us unpack this formula. Reading from left to right:

- $\arg\max_{a_t}$: choose the best current action.
- $\sum_{o_t, r_t}$: sum over all possible immediate responses.
- $\ldots \max_{a_m} \sum_{o_m, r_m}$: for all future time steps, maximize over actions and sum over environment responses. The alternation of $\max$ and $\sum$ reflects that future actions are chosen optimally while future observations are uncertain.
- $\sum_{k=t}^{m} r_k$: the total future reward.
- $\sum_{q} 2^{-|q|}$: the posterior weight on environment $q$, which is proportional to $2^{-|q|}$ restricted to programs consistent with what has been observed.

AIXI is simultaneously performing \textbf{Bayesian learning} (maintaining the posterior over environments via the Solomonoff prior) and \textbf{optimal planning} (choosing the reward-maximizing action given this posterior).
\end{definition}

\subsection{13.3.3 A Concrete Scenario}

\begin{example}[Cheese-in-a-maze]
An agent is placed in a maze. At each step, it chooses a direction (north, south, east, west). The environment returns the agent's new position (observation) and a reward of 1 if the agent reaches a cheese and 0 otherwise. The maze layout is fixed but unknown to the agent.

AIXI's posterior initially assigns weight to all computable maze-generating programs. Simple mazes (small, few walls) correspond to short programs and thus receive higher prior weight. As the agent explores and observes walls and open spaces, most maze programs become inconsistent with the data and are eliminated from the posterior. After sufficient exploration, the posterior concentrates on programs generating mazes similar to the true one. The agent then exploits this knowledge to navigate efficiently to the cheese.

The key insight is that AIXI's exploration is \textit{automatic}---it follows from the Bayesian framework. Actions that yield informative observations help distinguish between environments, which helps maximize \textit{future} reward. AIXI does not need a separate exploration strategy bolted on; exploration emerges from the expected reward maximization over the full horizon.
\end{example}

\subsection{13.3.4 Optimality}

\begin{theorem}[Pareto Optimality, Hutter 2002]
\textit{AIXI is Pareto-optimal over the class of all computable environments: there is no other policy $\pi$ that achieves higher expected reward than AIXI in some computable environment without achieving lower expected reward in another.}

This is the strongest optimality statement one can make without assuming which environment the agent is in. It says AIXI cannot be uniformly improved upon.
\end{theorem}

\begin{theorem}[Self-optimizing, Hutter 2002]
\textit{In any computable environment $\mu$, the per-step expected reward of AIXI converges to the optimal value achievable by a policy that knows $\mu$.}

In other words, AIXI eventually performs as well as an agent that knows the true environment. The "eventually" may take arbitrarily long for complex environments, but convergence is guaranteed.
\end{theorem}

\subsection{13.3.5 Limitations}

AIXI is a mathematical ideal, not a practical algorithm. Its limitations are severe and illuminating:

\begin{enumerate}
\item \textbf{Incomputability.} AIXI requires computing the Solomonoff prior and solving an infinite-depth expectimax search. Both are incomputable.
\end{enumerate}

\begin{enumerate}
\item \textbf{No self-reference.} AIXI models the environment as a program $q$ that is separate from the agent. It does not model itself as part of the environment. If the environment can inspect the agent's code or predict the agent's actions (as in game-theoretic settings), AIXI's model is incomplete. This is the \textbf{embedded agency} problem: real agents are \textit{within} their environments, not separate from them.
\end{enumerate}

\begin{enumerate}
\item \textbf{The grain of truth assumption.} AIXI assumes the true environment is computable. If the environment involves genuinely uncomputable elements (e.g., depends on a halting oracle), the Solomonoff prior assigns it zero weight.
\end{enumerate}

\begin{enumerate}
\item \textbf{Lack of self-modification.} AIXI's policy is defined by a fixed formula. It does not modify its own source code or architecture. It treats the problem as "fixed agent, unknown environment."
\end{enumerate}

\subsection{13.3.6 AIXItl: A Computable Approximation}

\begin{definition}[AIXItl, Hutter 2002]
Fix resource bounds: a maximum program length $l$ and a maximum computation time $t$ per step. The \textbf{AIXItl} agent is identical to AIXI except that the sum over programs is restricted to those of length $\leq l$ that run in time $\leq t$.

AIXItl is computable (since we only consider finitely many programs, each run for bounded time). It is not optimal in the same sense as AIXI, but it satisfies a different guarantee:
\end{definition}

\begin{theorem}
\textit{AIXItl is as good as any other agent of computational resources $(t, l)$. For any computable policy $\pi$ of description length $\leq l$ running in time $\leq t$, AIXItl achieves expected reward at least as high (up to terms that vanish with the horizon).}

This makes AIXItl the best possible use of bounded computational resources, in a well-defined sense. The practical difficulty is that the constants are astronomical: even modest values of $l$ yield an intractably large set of programs to consider.
\end{theorem}



\section{Gödel Machines}

\subsection{13.4.1 Motivation}

AIXI is a fixed agent reasoning about an unknown environment. But what about an agent that can improve \textit{itself}? A programmer who writes better versions of their own code? This is the problem of \textbf{self-modification}, and it leads directly to the Gödel Machine.

The key observation motivating the Gödel Machine is simple: if an agent can prove that a modification to its own code would increase expected future reward, then it should perform that modification. The challenge is making this idea rigorous.

\subsection{13.4.2 Structure}

\begin{definition}[Gödel Machine, Schmidhuber 2003]
A Gödel Machine is a computational system consisting of:

1. \textbf{A program $f$}, stored in memory and initially containing a proof searcher $s$. The program $f$ includes instructions for interacting with the environment and for searching for self-improvements.

2. \textbf{An axiom system $\mathcal{A}$}, formalized in a language powerful enough to describe:
   - The machine's own code and state.
   - The environment's properties (to whatever extent they are known).
   - Arithmetic and the theory of computation.

3. \textbf{A utility function $u$}, specifying the agent's objective (e.g., expected cumulative reward).

The machine operates as follows at each step:

- Execute the current instruction of $f$ (which may involve interacting with the environment, performing computations, or searching for proofs).
- The proof searcher $s$ systematically enumerates all possible proofs in $\mathcal{A}$.
- If $s$ finds a proof that replacing the current program $f$ with some $f'$ would yield strictly higher expected utility than continuing with $f$ (including continuing to search for better modifications), then \textbf{perform the rewrite}: replace $f$ with $f'$.
\end{definition}

\subsection{13.4.3 The Self-Referential Proof Search}

The proof searcher is the heart of the machine. It enumerates candidate theorems of the form:

> "Switching from the current policy $f$ to policy $f'$ at the current time, given the current state, yields higher expected utility than continuing with $f$ (which includes continuing the proof search)."

Formally, let $\text{state}(t)$ be the complete state of the machine at time $t$, and let $u(f, \text{state}(t))$ be the expected utility of running policy $f$ from state $\text{state}(t)$ onward. The proof searcher looks for a theorem of the form:

\[\mathcal{A} \vdash u(f', \text{state}(t)) > u(f, \text{state}(t))\]

where $f'$ is some candidate replacement program. If such a proof is found, the machine switches to $f'$.

The remarkable feature is that $f'$ is unrestricted: it can be \textit{any} program. In particular, $f'$ can contain a different proof searcher, a different utility function interpretation, a different environment model---anything. The only constraint is that the improvement must be \textit{proved}, not merely conjectured.

\subsection{13.4.4 Key Properties}

\textbf{Property 1: Global optimality guarantee.} If a beneficial rewrite exists and is provable within $\mathcal{A}$, the proof searcher will eventually find it (since it enumerates all proofs). The machine will not get stuck in a local optimum if the global optimum is reachable via a provable improvement.

\textbf{Property 2: Soundness of self-modification.} Every self-modification is backed by a formal proof. The machine never makes a change that it cannot justify. This prevents the pathological self-modifications that plague naive self-modifying systems (e.g., a program that "optimizes" itself by deleting all safety checks).

\textbf{Property 3: Self-referential stability.} The machine can rewrite the proof searcher itself. If there is a provably faster search strategy, it will be found and adopted. This makes the Gödel Machine optimally self-reflective in a precise sense: the proof searcher has no bias toward preserving itself.

\textbf{Property 4: Gödel's incompleteness applies.} By Gödel's first incompleteness theorem, any consistent axiom system $\mathcal{A}$ that is powerful enough to encode arithmetic contains true but unprovable statements. This means there may exist beneficial self-rewrites that the machine can never prove are beneficial. The Gödel Machine is not omniscient---it is limited to \textit{provably} beneficial modifications.

\begin{example}[Proof-driven self-improvement]
Suppose a Gödel Machine starts with a naive sorting algorithm (bubble sort, $O(n^2)$) as a subroutine used in its environment interactions. The axiom system includes the theory of arithmetic and facts about algorithm analysis. The proof searcher eventually constructs a proof that replacing bubble sort with merge sort ($O(n \log n)$) will decrease computation time for all inputs of size $\geq n_0$, thereby increasing expected utility (since faster computation means more time for other useful activities within a fixed time budget). Upon finding this proof, the machine rewrites its sorting subroutine.

Note that the machine does not merely test that merge sort is faster empirically. It \textit{proves} it, within the formal system, using the axioms of arithmetic and the definition of the algorithms. This proof may take enormous time to find, but once found, the improvement is guaranteed to be correct.
\end{example}

\subsection{13.4.5 Comparison with AIXI}

\begin{center}
\begin{tabular}{ccc}
\toprule
Property & AIXI & Gödel Machine \\
\midrule
Self-modification & No & Yes \\
Handles unknown environments & Yes (Bayesian) & Partially (via axioms) \\
Optimality guarantee & Pareto-optimal over computable envs & Optimal among provably good policies \\
Computability & Incomputable & Incomputable (in general) \\
Limitation principle & Halting problem & Gödel's incompleteness \\
\bottomrule
\end{tabular}
\end{center}

The two approaches address complementary aspects of the optimal agent problem. AIXI solves the learning problem (how to act in an unknown environment) but ignores self-modification. The Gödel Machine solves the self-modification problem (how to improve one's own code) but requires an axiom system that adequately describes the environment.



\section{OOPS: Optimal Ordered Problem Solver}

\subsection{13.5.1 Motivation}

Schmidhuber (2004) introduced the \textbf{Optimal Ordered Problem Solver (OOPS)} as a more concrete (though still idealized) approach to incremental self-improvement. The key idea is to reuse solutions to previously solved problems when solving new ones.

\subsection{13.5.2 Setup}

OOPS faces a sequence of problems $P_1, P_2, P_3, \ldots$ and maintains a growing \textit{code store} $C$. For each problem $P_k$:

\begin{enumerate}
\item OOPS performs \textbf{Levin search} (universal search, Levin 1973): enumerate all programs $p$ by increasing length, running each for time proportional to $2^{-|p|}$ times the total available time.
\item Crucially, programs in the search can \textit{call} previously found solutions as subroutines. If $P_3$ is similar to $P_1$, the solution to $P_1$ (already in the code store) can be invoked as a building block.
\item When a solution is found, it is added to $C$.
\end{enumerate}

\subsection{13.5.3 Bias-Optimality}

\begin{definition}[Bias-Optimality]
A search method is \textit{bias-optimal} if no other method with the same set of initial assumptions (prior over programs) can solve the given problem more than a constant factor faster.
\end{definition}

\begin{theorem}[Schmidhuber 2004]
\textit{OOPS is bias-optimal for its given problem sequence.}

This means OOPS is provably the fastest way (up to a constant) to solve a sequence of problems given Levin's universal prior, with the added benefit of reusing prior solutions.
\end{theorem}

\subsection{13.5.4 OOPS as a Dynamical System}

OOPS naturally generates a trajectory of increasingly capable solvers:

\[C_0, \quad C_1 = C_0 \cup \{ s_1\} , \quad C_2 = C_1 \cup \{ s_2\} , \quad \ldots\]

where $s_k$ is the solution to problem $P_k$. Each new solution potentially makes future searches faster (by providing callable subroutines). This is a dynamical system on the space of code stores, where the transition rule is "solve the next problem using all available code, then add the solution."

In our notation: $f_k = C_k$ is the agent's "program" at stage $k$, and $x_k = P_k$ is the current problem. The update is $f_{k+1} = \phi(f_k, x_k) = f_k \cup \{ \text{solve}(x_k, f_k)\} $. This is a concrete instance of the self-modifying dynamical system $(f, x) \mapsto (\phi(f,x), g(x))$ where the map itself grows monotonically.



\section{Synthesis: Self-Modifying Dynamical Systems}

We now step back and connect these constructions to the framework developed in previous chapters.

\subsection{13.6.1 Three Approaches to $(f, x) \mapsto (f', x')$}

Consider the general update rule where both the state and the map can change:

\[(f_t, x_t) \mapsto (f_{t+1}, x_{t+1}) = (\phi(f_t, x_t),\; f_t(x_t)).\]

The three constructions studied in this chapter correspond to three different choices of the meta-rule $\phi$:

\textbf{AIXI:} $\phi(f, x) = f$ (no self-modification). The map $f$ is the fixed AIXI policy. The state $x$ encodes the interaction history. AIXI updates its \textit{beliefs} about the environment, but this is a change in $x$ (data), not in $f$ (policy). The agent-environment system evolves as $(f, x) \mapsto (f, f(x))$---a standard iterated map with a fixed function.

\textbf{Gödel Machine:} $\phi(f, x) = f'$ where $f'$ is the result of a provably-beneficial rewrite if one is found, or $f$ otherwise. The meta-rule $\phi$ is constrained:

\[\phi(f, x) = \begin{cases} f' & \text{if } \mathcal{A} \vdash u(f', x) > u(f, x) \text{ is found at this step} \\ f & \text{otherwise.} \end{cases}\]

This is a self-modifying dynamical system with a \textit{conservative} transition rule: changes to $f$ are infrequent and always beneficial (relative to the axiom system).

\textbf{OOPS:} $\phi(f, x) = f \cup \{ \text{new solution}\} $. The meta-rule is \textit{monotone}: the code store only grows. This guarantees that capabilities never decrease, at the cost of ever-increasing memory usage.

\subsection{13.6.2 The Unconstrained Case}

Our general framework imposes no constraints on $\phi$. The meta-rule could be:

\begin{itemize}
\item \textbf{Arbitrary:} $\phi$ is any computable function. This allows studying phenomena like \textbf{model collapse} (where $f$ degenerates) or \textbf{creative transients} (long sequences of distinct $f$ values before settling).
\item \textbf{Stochastic:} $\phi$ selects $f'$ randomly from some distribution depending on $(f, x)$.
\item \textbf{Adversarial:} $\phi$ is chosen to make things as bad as possible (useful for studying robustness).
\end{itemize}

By not requiring provable improvement, we gain the ability to study the \textit{dynamics} of self-modification itself---the trajectories, attractors, transients, and bifurcations that arise when programs rewrite programs.

\begin{example}[Degenerate self-modification]
Let $\mathcal{F} = \{ f_0, f_1, f_2\} $ be three programs and let $\phi$ be defined by: $\phi(f_0, x) = f_1$, $\phi(f_1, x) = f_2$, $\phi(f_2, x) = f_2$ for all $x$. Then regardless of the initial state, the meta-dynamics converge to the fixed point $f_2$ in at most two steps. If $f_2$ is a degenerate program (e.g., one that always outputs the same action), this represents model collapse: a self-modifying system that destroys its own capabilities. A Gödel Machine would never reach $f_2$ (it could not prove that switching to a degenerate program improves utility), but the unconstrained dynamics allow it.
\end{example}

\subsection{13.6.3 What These Constructions Share}

Despite their differences, AIXI, Gödel Machines, and OOPS share several deep structural features:

\begin{enumerate}
\item \textbf{Universality.} Each construction ranges over \textit{all} programs (or all proofs, or all solutions), not a pre-specified model class. This universality is what gives the theoretical guarantees their power---and what makes the constructions incomputable.
\end{enumerate}

\begin{enumerate}
\item \textbf{Occam weighting.} Programs are weighted by $2^{-|p|}$, favoring short descriptions. This is the algorithmic formalization of Occam's razor and is the common thread linking Kolmogorov complexity (Chapter 8), the Solomonoff prior, and AIXI.
\end{enumerate}

\begin{enumerate}
\item \textbf{Incomputability barriers.} All three hit fundamental limits: the halting problem for Solomonoff/AIXI, Gödel's incompleteness for the Gödel Machine. These are not bugs but features---they delineate the boundary of what is achievable by any formal system.
\end{enumerate}



\section{Summary}

\begin{center}
\begin{tabular}{ccccc}
\toprule
Construction & Year & Core idea & Optimality & Limitation \\
\midrule
Solomonoff prior & 1964 & Universal mixture over programs & Dominates all computable measures & Incomputable \\
AIXI & 2000 & Bayesian RL with Solomonoff prior & Pareto-optimal & Incomputable, no self-modification \\
AIXItl & 2002 & AIXI with bounded resources & Optimal for given resources & Astronomical constants \\
Gödel Machine & 2003 & Provably beneficial self-rewrites & Optimal among provable policies & Gödel incompleteness \\
OOPS & 2004 & Incremental universal search with reuse & Bias-optimal & Idealized, large constants \\
\bottomrule
\end{tabular}
\end{center}

The progression from Solomonoff induction through AIXI to Gödel Machines represents a series of increasingly ambitious attempts to formalize the notion of an optimal intelligent agent. Solomonoff solves prediction. AIXI extends this to action. The Gödel Machine extends further to self-modification. Each step broadens the scope while preserving formal guarantees.

For the study of discrete dynamical systems on function spaces, these constructions provide the conceptual endpoints: AIXI is the best possible fixed map, and the Gödel Machine is the best possible self-modifying map (within the limits of provability). The space between them---self-modification without proof constraints---is where the interesting dynamics live.



\section{References}

\begin{itemize}
\item Hutter, M. (2005). \textit{Universal Artificial Intelligence: Sequential Decisions Based on Algorithmic Probability}. Springer.
\item Hutter, M. (2000). A theory of universal artificial intelligence based on algorithmic complexity. arXiv:cs/0004001.
\item Hutter, M. (2007). Universal algorithmic intelligence: A mathematical top-down approach. In \textit{Artificial General Intelligence}, 227--290.
\item Legg, S. (2008). \textit{Machine Super Intelligence}. PhD thesis, University of Lugano.
\item Levin, L. A. (1973). Universal sequential search problems. \textit{Problems of Information Transmission}, 9(3):265--266.
\item Schmidhuber, J. (2003). Gödel machines: Self-referential universal problem solvers making provably optimal self-improvements. arXiv:cs/0309048.
\item Schmidhuber, J. (2004). Optimal ordered problem solver. \textit{Machine Learning}, 54(3):211--254.
\item Schmidhuber, J. (2009). Ultimate cognition à la Gödel. \textit{Cognitive Computation}, 1(2):177--193.
\item Solomonoff, R. J. (1964). A formal theory of inductive inference. \textit{Information and Control}, 7(1):1--22 and 7(2):224--254.
\item Solomonoff, R. J. (1978). Complexity-based induction systems: Comparisons and convergence theorems. \textit{IEEE Transactions on Information Theory}, 24(4):422--432.
\item Zvonkin, A. K. and Levin, L. A. (1970). The complexity of finite objects and the development of the concepts of information and randomness by means of the theory of algorithms. \textit{Russian Mathematical Surveys}, 25(6):83--124.
\end{itemize}



\section{Recommended Reading}

For Solomonoff induction and AIXI:

\begin{itemize}
\item \textbf{Hutter (2005)} is the definitive monograph. The first half covers Solomonoff induction; the second half covers AIXI. Technically demanding but comprehensive.
\end{itemize}

\begin{itemize}
\item \textbf{Legg (2008)}, \textit{Machine Super Intelligence}, provides a more accessible introduction to the main ideas, including a formal definition of intelligence based on AIXI.
\end{itemize}

\begin{itemize}
\item \textbf{Rathmanner, S. and Hutter, M.} (2011). "A philosophical treatise of universal induction." \textit{Entropy}, 13(6):1076--1136. Covers the philosophical foundations and addresses common objections.
\end{itemize}

For Gödel Machines:

\begin{itemize}
\item \textbf{Schmidhuber (2003)} is the original paper and remains the best introduction. Available on arXiv as cs/0309048.
\end{itemize}

\begin{itemize}
\item \textbf{Schmidhuber (2009)}, "Ultimate cognition à la Gödel," provides a shorter and more accessible overview.
\end{itemize}

For OOPS and incremental self-improvement:

\begin{itemize}
\item \textbf{Schmidhuber (2004)}, "Optimal ordered problem solver," in \textit{Machine Learning}. The original paper, with full technical details.
\end{itemize}

For broader context on AGI theory:

\begin{itemize}
\item \textbf{Goertzel, B. and Pennachin, C.} (eds.) (2007). \textit{Artificial General Intelligence}. Springer. Collection of papers from the AGI community, including Hutter's overview of AIXI.
\end{itemize}

For embedded agency and limitations of AIXI:

\begin{itemize}
\item \textbf{Demski, A. and Garrabrant, S.} (2019). "Embedded Agency." MIRI Technical Report. Analyzes the assumptions underlying AIXI and identifies problems that arise when the agent is embedded within its environment.
\end{itemize}

For practical approximations:

\begin{itemize}
\item \textbf{Veness, J. et al.} (2011). "A Monte-Carlo AIXI Approximation." \textit{Journal of Artificial Intelligence Research}, 40:95--142. The most successful practical approximation to AIXI, with experiments on Atari games.
\end{itemize}
